\appendix
\section{Establishing Equation \protect\ref*{provedintheappendix}}

\label{app}

%
\subsection{A clean calculation establishing a loose version of equation
\protect\ref*{provedintheappendix}}

For clarity of exposition, we begin by presenting a relatively clean
calculation that establishes a slightly loose version of 
\Cref{provedintheappendix}. Using just this looser bound, we would be able to
establish that \Cref{provedintheappendix} holds (with the constant $1/2$
replaced by a slightly smaller constant) so long as we set $d_1$ to be $
\Theta\left(w^{1/3}/\log w \right)$. A slightly more involved calculation
(cf. \Cref{tightsec}) is required to establish 
\Cref{provedintheappendix} for our desired value of $d_1= \lfloor (w/c)^{1/3}
\rfloor$.

Expression \eqref{aboveexp} equals
\begin{align}
\left(\frac{w}{ci^2}\right)^{d_1-1} \cdot \frac{i^2}{\left((d_1)!\right)^2}
\cdot \prod_{j=1, j \neq i}^{d_1}\left( \left|j-i| \cdot |j+i\right| - \frac{
c i^2 j^2}{w}\right)  \notag \\
= \left(\frac{w}{ci^2}\right)^{d_1-1} \cdot \frac{i^2}{\left((d_1)!\right)^2}
\cdot \prod_{j=1, j \neq i}^{d_1}\left( \left|j-i| \cdot |j+i\right| \right)
\cdot \left(1-\frac{c i^2 j^2}{w \cdot |j-i| |j+i|}\right)  \notag \\
= \left(\frac{w}{ci^2}\right)^{d_1-1} \cdot \frac{(d_1+i)! (d_1-i)!}{
2\left((d_1)!\right)^2} \cdot \prod_{j=1, j \neq i}^{d_1} \left(1-\frac{c
i^2 j^2}{w \cdot |j-i| |j+i|}\right)  \label{looseversion} \\
\geq \left(\frac{w}{ci^2}\right)^{d_1-1} \cdot \frac{1}{2} \cdot \prod_{j=1,
j \neq i}^{d_1} \left(1-\frac{c i^2 j^2}{w \cdot |j-i| |j+i|}\right)  \notag
\\
\geq \left(\frac{w}{ci^2}\right)^{d_1-1} \cdot \frac{1}{2} \cdot
\left(1-\sum_{j=1, j \neq i }^{d_1}\frac{c i^2 j^2}{w \cdot |j-i| |j+i|}
\right)  \notag \\
\geq \left(\frac{w}{ci^2}\right)^{d_1-1} \cdot \frac{1}{2} \cdot \left(1-
\frac{c i^2}{w}\sum_{j=1, j \neq i }^{d_1}\frac{j^2}{ |j-i| |j+i|}\right).
\label{loosemiddle}
\end{align}

Let us consider the expression $\sum_{j=1, j \neq i }^{d_1}\frac{j^2}{ |j-i|
|j+i|}$. If $i^2 \not \in [j^2/2, 3j^2/2]$, then the $j$'th term in this sum
is at most $2$. Hence, letting $H_i$ denote the $i$th Harmonic number and
using the fact that $H_i \leq \ln(i+1)$,
\begin{align}
& \sum_{j=1, j \neq i }^{d_1}\frac{j^2}{ |j-i| |j+i|}  \notag \\
& \leq 2 \cdot d_1 + \sum_{j = \lfloor \sqrt{2/3} i \rfloor}^{\lfloor \sqrt{2
} i \rfloor} \frac{j^2}{|j-i| |j+i|}  \notag \\
& \leq 2 \cdot d_1 + \sum_{j = \lfloor \sqrt{2/3} \cdot i \rfloor}^{\lceil
\sqrt{2} \cdot i \rceil} \frac{j}{|j-i|}  \notag \\
& \leq 2 d_1 + \sqrt{2} \cdot i \cdot \sum_{j=1}^{(\sqrt{2}-1) \cdot i} 2/j
\notag \\
& \leq 2 d_1 + 2 \sqrt{2} \cdot i \cdot H_{i} \leq 2d_1 + 2 \sqrt{2} i
\ln(i+1).  \label{eq:ieq}
\end{align}

We conclude that if $d_1$ were set to a value less than $w^{1/3}/(100 \cdot
c^2 \cdot \ln(w))$ (rather than to $\lfloor \left(w/c\right)^{1/3}\rfloor$),
then Expression \eqref{loosemiddle} is at least
\begin{equation}
\left(\frac{w}{ci^2}\right)^{d_1-1} \cdot \frac{1-1/c}{2}.
\end{equation}

%
\subsection{The tight bound}

\label{tightsec} To obtain the tight bound, we need a tighter sequence of
inequalities following Expression \eqref{looseversion}. Specifically,
Expression \eqref{looseversion} is bounded below by:

\begin{align}
\geq \left(\frac{w}{ci^2}\right)^{d_1-1} \cdot \frac{1}{2} \left(1+ \frac{i}{
2d_1}\right)^i \cdot \prod_{j=1, j \neq i}^{d_1} \left(1-\frac{c i^2 j^2}{w
\cdot |j-i| |j+i|}\right)  \notag \\
\geq \left(\frac{w}{ci^2}\right)^{d_1-1} \cdot \frac{1}{2} \cdot
e^{i^2/(2d_1)} \cdot \prod_{j=1, j \neq i}^{d_1} \left(1-\frac{c i^2 j^2}{w
\cdot |j-i| |j+i|}\right)  \notag \\
\geq \left(\frac{w}{ci^2}\right)^{d_1-1} \cdot \frac{1}{2} \cdot
e^{i^2/(2d_1)} \cdot \prod_{j=1, j \neq i}^{d_1} \left(1-\frac{c i^2 j^2}{w
\cdot |j-i| |j+i|}\right)  \label{tightbegin}
\end{align}

The rough idea of how to proceed is as follows. \Cref{eq:ieq}
implies that for $i \ll w^{1/3}/\ln w$, the factor
\begin{equation*}
F_1 := \prod_{j=1, j \neq i}^{d_1} \left(1-\frac{c i^2 j^2}{w \cdot |j-i|
|j+i|}\right)
\end{equation*}
is at some a positive constant, and hence Expression \eqref{tightbegin} is
bounded below by the desired quantity. If $i \gtrsim w^{1/3}/\ln w$, then
\Cref{eq:ieq} does not yield a good bound on this factor, leaving
open the possibility that this factor is subconstant. But in this case, the
factor $F_2:= e^{i^2/(2d_1)} \geq e^{\tilde{\Omega}(d_1)}$, and the
largeness of $F_2$ dominates the smallness of $F_1$.

In more detail, let $x_{i, j} = \frac{ci^2j^2}{w \cdot |i-j| j+i||}$. Then
for all $i \neq j$ such that $i, j \leq d_1$,
\begin{equation}  \label{xijbound}
x_{i,j} \leq \frac{c \cdot d_1^2 (d_1-1)^2}{(2d_1-1) \cdot w} \leq \frac{c
\cdot d_1^3}{2w} \leq 1/2,
\end{equation}
where in the final inequality we used the fact that $d_1 \leq (w/c)^{1/3}$.

Using the fact that $1-x \geq e^{-x-x^2}$ for all $x \in [0, 1/2]$, we can
write
\begin{equation*}
F_1 \geq \prod_{j=1, j\neq i}^{d_1} e^{-x_{i,j} - x_{i, j}^2}.
\end{equation*}
Hence,
\begin{align*}
F_1 \cdot F_2 \geq \exp\left(i^2/(2d_1) - \sum_{j=1, j\neq i}^{d_1} -x_{i,j}
- x_{i, j}^2\right).
\end{align*}

From Equations \eqref{eq:ieq} and \eqref{xijbound}, we know that
\begin{align*}
\sum_{j=1, j\neq i}^{d_1} x_{i,j} + x_{i, j}^2 \leq \frac{c i^2}{w} \cdot
\left(3d_1 + 3 \sqrt{2} i \ln(i+1)\right) \leq \frac{ci^2}{w} \cdot \left(4
d_1 \ln(d_1)\right).
\end{align*}
Hence,
\begin{align*}
F_1 \cdot F_2 & \geq \exp\left(i^2/(2d_1) - \frac{ci^2}{w} \cdot 4 c
\ln(d_1)\right) \\
&= \exp\left(i^2 \left(\frac{1}{2d_1} - \frac{4c^2 \ln(d_1)}{w}\right)\right)
\\
& \geq \exp\left(i^2 \cdot \frac{1}{2d_1} \cdot (1-o(1))\right) \\
& \geq 1.
\end{align*}

\Cref{provedintheappendix} follows.



%
